% ═══════════════════════════════════════════════════════════════════════════════
%  CAPÍTULO — State
% ═══════════════════════════════════════════════════════════════════════════════

\chapter{State}

% ─── Situación Problema ──────────────────────────────────────────────────────
\section{Situación Problema}

Una reserva en Airbnb atraviesa varios estados a lo largo de su ciclo
de vida: \emph{solicitada}, \emph{confirmada}, \emph{en curso},
\emph{completada} y \emph{cancelada}. El comportamiento permitido varía
drásticamente según el estado: una reserva solicitada puede ser rechazada
por el anfitrión, pero una completada no; se puede cancelar una reserva
confirmada, pero con penalizaciones distintas según la antelación; no
es posible dejar reseña hasta que la reserva esté completada.

Implementar estas reglas con condicionales sobre un atributo de estado
produce un código lleno de \texttt{if-else} anidados que crece con cada
nuevo estado o transición. El riesgo de omitir una combinación inválida
es alto, y el código resulta difícil de leer y extender.

% ─── Solución ────────────────────────────────────────────────────────────────
\section{Solución}

El patrón State representa cada estado del ciclo de vida como una clase
concreta que implementa la interfaz \texttt{EstadoReserva}. La clase
\texttt{Reserva} delega la respuesta a cada acción —\texttt{confirmar()},
\texttt{cancelar()}, \texttt{completar()}, \texttt{dejarResena()}— al
objeto de estado que tiene asignado en ese momento. Cada estado concreto
sabe exactamente qué transiciones son válidas, qué lógica ejecutar y a
qué estado siguiente moverse.

El código de cada estado es compacto y autocontenido. Añadir un nuevo
estado o modificar las reglas de transición de uno existente no requiere
alterar los demás. La clase \texttt{Reserva} queda libre de condicionales
y concentrada en su responsabilidad principal.

% ─── Diagrama de clases ─────────────────────────────────────────────────────
\begin{figure}[H]
    \centering
    % \includegraphics[width=\textwidth]{imagenes/abstract_factory_diagrama.png}
    \caption{Diagrama de clases del patrón Abstract Factory}
\end{figure}


% ─── Roles de las Clases ────────────────────────────────────────────────────
\section{Roles de las Clases}

% Explicar la responsabilidad de cada clase/interfaz

\begin{itemize}
    \item \textbf{AbstractFactory:}
    \item \textbf{ConcreteFactory:}
    \item \textbf{AbstractProduct:}
    \item \textbf{ConcreteProduct:}
    \item \textbf{Client:}
\end{itemize}


% ─── Main ───────────────────────────────────────────────────────────────────
\section{Main}

% Explicar qué realiza el programa principal

\begin{lstlisting}[language=Java, caption=NombreClase]
 
\end{lstlisting}


% ─── Salida del Main ────────────────────────────────────────────────────────
\section{Salida del Main}

\begin{lstlisting}[language=Java, caption=NombreClase]
 
\end{lstlisting}


% ─── Clases ─────────────────────────────────────────────────────────────────
\section{Clases}

% ─── Clase 1 ────────────────────────────────────────────────────────────────
\subsection{NombreClase}

\begin{lstlisting}[language=Java, caption=NombreClase]
 
\end{lstlisting}


% ─── Clase 2 ────────────────────────────────────────────────────────────────
\subsection{NombreClase}

\begin{lstlisting}[language=Java, caption=NombreClase]
 
\end{lstlisting}


% ─── Clase 3 ────────────────────────────────────────────────────────────────
\subsection{NombreClase}

\begin{lstlisting}[language=Java, caption=NombreClase]
 
\end{lstlisting}


% ─── Conclusiones ───────────────────────────────────────────────────────────
\section{Conclusiones}

% Explicar:
% - Ventajas del patrón
% - Cuándo usarlo
% - Beneficios obtenidos
% - Posibles desventajas
% - Relación con principios SOLID